\chapter{Bảng Kịch Bản Kiểm Thử Hệ Thống Chi Tiết}
\label{app:testcases}

Phụ lục này trình bày chi tiết các kịch bản kiểm thử (Test Cases) được thực hiện trong quá trình tích hợp hệ thống ở Tuần 11 và Tuần 12, nhằm thẩm định tính năng và độ ổn định của thiết bị nhúng.

\section{Kịch bản kiểm thử tích hợp phần cứng và ngoại vi}
Danh sách các bài test kiểm tra hoạt động cơ bản của vi điều khiển ESP32 và các module phần cứng được mô tả trong Bảng \ref{tab:testsuitehardware}.

\begin{longtable}{llp{5.5cm}p{5.5cm}}
\caption{Kịch bản kiểm thử phần cứng và ngoại vi cục bộ} \label{tab:testsuitehardware} \\
\toprule
\textbf{Mã Test} & \textbf{Chức năng} & \textbf{Quy trình kiểm thử} & \textbf{Kết quả thực tế thu được} \\
\midrule
\endfirsthead
\toprule
\textbf{Mã Test} & \textbf{Chức năng} & \textbf{Quy trình kiểm thử} & \textbf{Kết quả thực tế thu được} \\
\midrule
\endhead
\bottomrule
\endfoot
\bottomrule
\endlastfoot

\textbf{TC-HW-01} & Khối nguồn & 
1. Cắm nguồn DC 5V. \newline
2. Đo điện áp tại ngõ ra nguồn bảo vệ. \newline
3. Đo điện áp tại chân 3V3 của ESP32. &
\begin{itemize}
    \item Điện áp ngõ ra nguồn ổn định ở $5.02$ V.
    \item Chân 3V3 của ESP32 đạt $3.29$ V. Không xảy ra sụt áp khi kết nối Wi-Fi.
\end{itemize} \\
\midrule

\textbf{TC-HW-02} & Mạch quạt hút và Relay &
1. ESP32 khởi động kéo chân \texttt{RELAY\_PIN} lên HIGH. \newline
2. Đo điện áp kích chân của Transistor đệm. &
\begin{itemize}
    \item Relay đóng phát tiếng tách giòn, quạt hút lập tức quay mạnh ổn định.
    \item Transistor đệm hoạt động mát, không quá nhiệt.
\end{itemize} \\
\midrule

\textbf{TC-HW-03} & Màn hình OLED SSD1306 &
1. Khởi động thiết bị. \newline
2. Đọc địa chỉ I2C quét bus. &
\begin{itemize}
    \item Phát hiện thiết bị I2C tại địa chỉ \texttt{0x3C}.
    \item OLED hiển thị Splash Screen đúng hình ảnh thiết lập.
\end{itemize} \\
\midrule

\textbf{TC-HW-04} & Cảnh báo Led \& Buzzer &
1. Ghi đè điều kiện logic \texttt{warning = true} trong code. \newline
2. Đo chu kỳ nháy chân còi buzzer. &
\begin{itemize}
    \item Đèn LED đỏ nhấp nháy chu kỳ 0.5s.
    \item Còi buzzer phát tiếng kêu tít tít ngắt quãng đúng tần số.
\end{itemize} \\
\bottomrule
\end{longtable}

\newpage
\section{Kịch bản kiểm thử cơ chế truyền thông và xử lý lỗi mạng}
Danh sách các bài test kiểm tra khả năng truyền dữ liệu đám mây, xử lý lỗi và lưu trữ offline cache được mô tả trong Bảng \ref{tab:testsuitenetwork}.

\begin{longtable}{llp{5.5cm}p{5.5cm}}
\caption{Kịch bản kiểm thử truyền thông và xử lý lỗi mạng} \label{tab:testsuitenetwork} \\
\toprule
\textbf{Mã Test} & \textbf{Chức năng} & \textbf{Quy trình kiểm thử} & \textbf{Kết quả thực tế thu được} \\
\midrule
\endfirsthead
\toprule
\textbf{Mã Test} & \textbf{Chức năng} & \textbf{Quy trình kiểm thử} & \textbf{Kết quả thực tế thu được} \\
\midrule
\endhead
\bottomrule
\endfoot
\bottomrule
\endlastfoot

\textbf{TC-NET-01} & Gửi HTTP POST lên Supabase &
1. Thiết bị kết nối Wi-Fi ổn định. \newline
2. Thực hiện đo đạc và gọi hàm POST. \newline
3. Truy cập bảng trên Supabase Console. &
\begin{itemize}
    \item Log Serial báo HTTP Code: 201 Created.
    \item Dữ liệu chèn thành công vào bảng với đầy đủ trường nhiệt độ, bụi, khí, trạng thái.
\end{itemize} \\
\midrule

\textbf{TC-NET-02} & Xử lý lỗi cảm biến trả về NaN &
1. Rút giắc cắm dữ liệu cảm biến DHT22 (giả lập hỏng cảm biến). \newline
2. Theo dõi log JSON và dữ liệu gửi lên DB. &
\begin{itemize}
    \item ESP32 nhận biết giá trị đọc là \texttt{NAN}.
    \item Log payload gửi đi: \texttt{"temperature": null}.
    \item Bảng trên DB chèn giá trị \texttt{null} thành công, không gây lỗi cú pháp JSON.
\end{itemize} \\
\midrule

\textbf{TC-NET-03} & Mất Wi-Fi đột ngột (Offline) &
1. Thiết bị đang đo đạc. \newline
2. Tắt bộ phát Wi-Fi access point. \newline
3. Theo dõi hành vi ghi file trong SPIFFS. &
\begin{itemize}
    \item TaskNetwork phát hiện mất kết nối Wi-Fi.
    \item Gọi hàm cache của DataManager.
    \item Log báo: \texttt{[DataManager] Record cached offline} xuống file \texttt{/cache.txt} thành công. Màn hình OLED hiển thị "Wi-Fi: Disconnect".
\end{itemize} \\
\midrule

\textbf{TC-NET-04} & Đồng bộ bù khi có mạng trở lại &
1. Bật lại Wi-Fi access point. \newline
2. Chờ chu kỳ đo tiếp theo của TaskSensors. &
\begin{itemize}
    \item TaskNetwork tự động reconnect thành công.
    \item Quét và phát hiện \texttt{/cache.txt}.
    \item Đọc từng dòng, chèn vào DB với trạng thái \texttt{status = 3} (cached offline).
    \item Log báo: \texttt{Cached record successfully posted}. File cache gốc tự động xóa sạch.
\end{itemize} \\
\bottomrule
\end{longtable}
